\chapter{CONCLUSIONES}
% Resumen de lo conseguido en el trabajo
En este documento se ha profundizado en el desarrollo de aplicaciones desde el punto de vista de la Ingeniería Biomédica para dar solución a un problema real, la localización de personas desaparecidas. Se ha diseñado una herramienta que trabaja con gran cantidad de datos con el objetivo de ayudar a salvar vidas. En concreto, a lo largo del trabajo se ha diseñado y desarrollado un prototipo de aplicación multiplataforma orientada al apoyo en la búsqueda y rescate de personas desaparecidas. 

% ¿Cómo la investigación aporta luz sobre el problema?
La investigación da como resultado una solución que permite digitalizar y agilizar flujos de comunicación y coordinación que actualmente son mayoritariamente manuales o fragmentados, para ello combina alertas, geolocalización, gestión estructurada de incidencias y mecanismos de colaboración entre grupos de rescate y ciudadanía. Además, al estudiar tanto el contexto social y operativo de las desapariciones como el fundamento biológico de la orientación humana, el trabajo conecta los conocimientos teóricos con las funcionalidades propuestas, ofreciendo un marco donde reducir tiempos de respuesta y mejorar la toma de decisiones en situaciones de emergencia.

% Recapitulación de objetivos y discusión de porqué no se lograron algunos
El objetivo principal del trabajo que era el de diseñar una aplicación de apoyo en la búsqueda y rescate de personas desaparecidas, se puede decir que se ha llevado a cabo con éxito puesto que de forma general se han definido las funcionalidades que debería tener la aplicación ideal. Respecto a los subobjetivos fijados, se completaron aquellos relativos al estudio del contexto actual y del origen biológico de la orientación, el análisis de herramientas existentes y la selección de funcionalidades relevantes, así como la definición general de la aplicación y las características del prototipo. 

Se realizó un estudio exhaustivo del contexto actual en torno a las personas desaparecidas mediante revisión bibliográfica y se profundizó en el origen biológico de la orientación, estudiando el concepto de ``GPS cerebral'' y su relación con la conducta humana de orientación. A continuación se investigaron las herramientas actualmente empleadas en búsquedas y rescates, tanto específicas como adaptadas, y se llevó a cabo un análisis comparativo que sirvió para seleccionar las funcionalidades más relevantes a incorporar en el diseño. Sobre esa base teórica se ha definido la aplicación a nivel conceptual y se han definido requisitos, casos de uso, arquitectura y pantallas principales para guiar el desarrollo de los primeros prototipos.

En lo relativo al prototipo, en el trabajo se ha desarrollado un prototipo básico funcional que integra frontend y backend utilizando Visual Studio Code, React Native y MariaDB con pantallas, procesos básicos de alertas y gestión de inicidencias. Sin embargo, por la complejidad técnica del proyecto y las limitaciones temporales del proyecto, determinadas funcionalidades del prototipo no se han llegado a implementar completamente, como el procedimiento de inicio de sesión seguro, los mapas o la geolocalización de participantes. 

Asimismo, se han sentado y probado las bases conceptuales de un sistema híbrido de predicción de localización de la persona desaparecida basado en características de la persona desaparecida y la zona demográfica de desaparición, esta base no pudo ser estudiada en profundidad ni validada con datos reales porque no existen datos abiertos accesibles ni posibilidad de obtener datos policiales anonimizados, ni publicaciones previas. Sin embargo, la experimentación realizada permitió realizar la validación del flujo del sistema predictivo propuesto aunque no la validación completa de los resultados predictivos del modelo.

% Reflexiones personales.
Personalmente, la complejidad de este proyecto me ha exigido un alto grado de investigación, puesta en común de distintos procedimientos actuales y abordaje del problema desde una perspectiva interdisciplinar poco habitual. Por otro lado, el desarrollo del prototipo me ha exigido un aprendizaje profundo en tecnologías full stack y en herramientas concretas como React Native, Expo, Visual Studio Code, MariaDB, JavaScript y Node.js... Además, esta idea ha resultado ser mucho más ambiciosa de lo esperado, especialmente en lo relativo al desarrollo del prototipo y el acceso a los datos, y está claro que avanzar hacia una versión más operativa exige de más tiempo para la implementación del prototipo y acuerdos institucionales para profundizar en el estudio de la posibilidad de predicción de la localización de una persona en base a su información personal.


% Implicaciones de los hallazgos en el contexto.
Con todo, se ha confirmado que una herramienta de estas características donde exista automatización de alertas y gestión digital de la información puede transformar la coordinación operativa en casos de personas desaparecidas al optimizar la asignación de recursos, mejorar la trazabilidad de acciones y aumentar las probabilidades de éxito en búsquedas críticas. Asimismo, la aplicación propuesta puede servir como plataforma de apoyo para profesionales y colectivos voluntarios, fomentar la colaboración interinstitucional y potenciar iniciativas internacionales dada la naturaleza global del problema. Sin embargo, la adopción efectiva en España depende de acuerdos formales con las fuerzas y cuerpos de seguridad para asegurar los protocolos. Sumado a esto, la integración con registros oficiales y sistemas como el PDyRH podría convertir esta aplicación en un recurso de alto valor operativo y social donde se garantiza una interoperabilidad con sistemas nacionales existentes.
